\clearpage
\phantomsection% 
\addcontentsline{toc}{chapter}{RINGKASAN}
\pagestyle{alternatingstyle}
\thispagestyle{alternatingstyle}

\begin{center}
	\large \bfseries \MakeUppercase{Ringkasan}\\
	\normalsize \normalfont {\thetitle}\\
	\normalsize \normalfont {\theauthor}\\
	\bigskip
	
	\normalsize \normalfont \justifying
	\normalsize \normalfont \justifying
	 Media sosial mendorong meningkatnya penggunaan meme sebagai bentuk komunikasi digital multimodal yang menggabungkan elemen visual dan teks. Meme tidak hanya berfungsi sebagai hiburan, tetapi juga dapat menyampaikan pesan implisit yang kompleks, termasuk konten sensitif seperti \textit{self-harm}. Konten seperti ini berisiko menimbulkan dampak negatif, terutama ketika disajikan secara tidak langsung melalui humor, simbol, atau permainan kata. Dalam Tugas Akhir ini, satu meme dapat terlihat netral secara visual, tetapi menyimpan makna yang berbahaya jika dibaca bersama dengan teksnya. Di sisi lain, penelitian deteksi meme \textit{self-harm} masih terbatas karena tantangan memahami konteks implisit, ambiguitas makna, serta keterbatasan ketersediaan data yang benar-benar representatif. Oleh karena itu, kebutuhan akan pendekatan yang mampu memanfaatkan informasi lintas modalitas atau kontekstual menjadi semakin penting dalam penelitian ini.
	 
	 Berdasarkan hal tersebut, penelitian ini merumuskan tiga fokus utama, yaitu pembangunan model klasifikasi biner multimodal berbasis \textit{Contrastive Language-Image Pre-training} (CLIP) dan \textit{Efficiently Learning an Encoder that Classifies Token Replacements Accurately} (ELECTRA) dengan pendekatan \textit{intermediate fusion}, analisis perbandingan pendekatan multimodal dan unimodal, serta evaluasi kinerja model menggunakan \textit{confusion matrix}. Tujuan penelitian adalah mengembangkan model deteksi meme \textit{self-harm}, mengukur efektivitas integrasi representasi visual dan tekstual, serta mengevaluasi performa model melalui \textit{accuracy}, \textit{precision}, \textit{recall}, dan \textit{F1-score}.
	 
	 Metodologi penelitian dimulai dari pengumpulan data primer dan sekunder yang berasal dari media sosial dan \textit{dataset} publik. Selanjutnya dilakukan rekayasa \textit{dataset} sintetis untuk memperkuat representasi kelas \textit{self-harm}. Data yang digunakan terdiri atas 524 meme \textit{self-harm} sintetis dan 1.500 meme non \textit{self-harm} dari \textit{dataset} publik. Seluruh data dianotasi menggunakan \textit{annotation guideline} berbasis literatur agar proses pelabelan lebih konsisten. Untuk mendukung anotasi, digunakan pendekatan \textit{ensemble pseudo-labeling} berbasis \textit{Large Language Model}, sehingga label yang dihasilkan dapat dipadukan dengan peninjauan manual. Pendekatan ini membantu menjaga konsistensi label pada data yang rentan memiliki ambiguitas semantik. Pengembangan model dilakukan dalam dua skema, yaitu unimodal dan multimodal, dengan multimodal menggabungkan representasi visual dari CLIP dan representasi tekstual dari ELECTRA melalui \textit{intermediate fusion}. Pendekatan ini dipilih agar hubungan semantik lintas modalitas dapat dipelajari secara lebih efektif, terutama pada meme yang maknanya dibangun dari kombinasi gambar dan teks.
	 
	 Pada tahap evaluasi, hasil menunjukkan bahwa pendekatan multimodal memiliki performa yang lebih baik dibandingkan pendekatan unimodal. Model multimodal memperoleh \textit{F1-score} sebesar 96,1\%, sedangkan model unimodal gambar dan unimodal teks masing-masing memperoleh \textit{F1-score} sebesar 91,5\% dan 92,0\%. Hal ini menunjukkan bahwa integrasi fitur visual dan tekstual membantu model memahami konteks secara lebih utuh dan mengurangi kesalahan klasifikasi. Hasil ini menegaskan bahwa informasi visual dan teks pada meme saling melengkapi, terutama saat salah satu modalitas tidak cukup untuk menangkap makna secara langsung. Dengan demikian, pendekatan multimodal lebih sesuai untuk mendeteksi meme \textit{self-harm} yang sering kali bersifat implisit. Hal ini menunjukkan bahwa representasi multimodal lebih mampu menangani variasi ekspresi, gaya bahasa tidak langsung, dan konteks visual yang berubah-ubah.

	Selain evaluasi model, dilakukan analisis perbandingan antara evaluasi model dan annotator manusia pada 19 data uji yang mengandung ambiugitas dan konteks kompleks. Hasil evaluasi kualitas anotasi menunjukkan 14 data berada pada kategori \textit{agreement} dan 5 data berada pada kategori \textit{disagreement}. Perbedaan ini muncul terutama pada meme yang ambigu, seperti meme yang mengandung sarkasme, humor gelap, atau permainan kata, sehingga makna \textit{self-harm} tidak selalu dapat dikenali secara langsung. Dalam beberapa kasus, visual dan teks justru memberi sinyal yang berlawanan sehingga interpretasi menjadi semakin sulit. Mengatasi perbedaan tersebut, dilakukan peninjauan ulang secara manual berdasarkan \textit{annotation guideline} untuk menetapkan label \textit{ground truth} yang konsisten. Proses ini menunjukkan bahwa evaluasi manusia tetap diperlukan sebagai mekanisme kontrol kualitas, khususnya pada data yang memiliki risiko salah tafsir tinggi.

	Evaluasi model pada data uji tersebut menghasilkan \textit{accuracy} 63,2\%, \textit{precision} 88,9\%, \textit{recall} 57,1\%, dan \textit{F1-score} 69,6\%. Nilai tersebut menunjukkan bahwa pada skenario evaluasi model dan manusia, performa model masih relatif rendah. Tingginya nilai \textit{precision} menunjukkan bahwa prediksi \textit{self-harm} yang dihasilkan model cenderung tepat, tetapi nilai \textit{recall} yang lebih rendah menandakan bahwa masih ada sejumlah kasus \textit{self-harm} yang terlewat atau diklasifikasikan sebagai non \textit{self-harm}. Hal ini mengindikasikan adanya \textit{false negative}, terutama pada meme dengan konflik antara teks dan visual. Di sisi lain, adanya \textit{false positive} menunjukkan bahwa model masih kesulitan membedakan antara ekspresi emosional umum, humor, dan indikasi \textit{self-harm}. Hal ini menegaskan bahwa beberapa contoh berada di luar pola data yang dipelajari model, sehingga termasuk dalam kategori data di luar distribusi (\textit{out-of-distribution}).

	Secara keseluruhan, penelitian ini menunjukkan bahwa model multimodal CLIP dan ELECTRA dengan pendekatan \textit{intermediate fusion} efektif dan lebih unggul dibandingkan pendekatan unimodal untuk mendeteksi meme \textit{self-harm}. Integrasi dua modalitas memberi kontribusi penting dalam memperkaya pemahaman konteks, terutama ketika pesan tersirat tidak tercermin secara eksplisit pada salah satu modalitas saja. Namun, hasil evaluasi model dan manusia memperlihatkan bahwa performa masih memiliki keterbatasan pada kasus ambigu dan data \textit{out-of-distribution}. Oleh karena itu, penelitian selanjutnya disarankan melibatkan ahli psikologi atau kesehatan mental dalam proses anotasi agar validitas dan sensitivitas label meningkat. Selain itu, penambahan variasi \textit{dataset} dengan konteks yang lebih kompleks, seperti sarkasme, humor gelap, dan variasi ekspresi linguistik lain, diperlukan agar model lebih robust. Secara praktis, model ini dapat dikembangkan sebagai alat bantu skrining awal, bukan pengganti penilaian manusia. Pengembangan sistem deteksi \textit{real-time} juga dapat menjadi arah lanjutan, dengan tetap mempertimbangkan aspek etika, privasi, dan dampak sosial dari penggunaan sistem tersebut.

	\vfill
	
\end{center}
\clearpage